\section{结论}

本文完成了车载平衡滚球系统的硬件、视觉、通信与控制设计。车辆采用八路光电循迹PD外环和双轮PI速度内环，
K230完成钢球检测、预测选通和位置滤波，STM32通过位置PID、球速阻尼及车辆运动前馈驱动PPR摆杆；
系统同时具备通信超时、视觉丢球、位置越界、输出限幅和渐进回中保护。

一次保留完整日志的偏置目标整车试验用时\SI{28.00}{s}，平均绝对误差\SI{0.78}{cm}，最大绝对误差\SI{1.76}{cm}，
全程未触发丢球保护。由此可确认系统已实现整车闭环并满足\SI{30}{s}时间门槛，但尚未满足全过程\SI{1}{cm}精度要求，
主要超限集中在右转阶段。R1图传具备实验性软件实现，R2至R5核心功能已在代码中完成；由于未保留相应正式重复测试资料，
本文不填写无法复核的成绩。该结论既反映了当前系统实际完成度，也为后续围绕弯道前馈、机械间隙和试验留档的改进提供了明确依据。
